\subsection{DIFFERENTIATION OF COMPARISON RELATIONS}

\pandocbounded{\includegraphics[keepaspectratio]{images/part002_390c97e1.jpg}}

Propositions 1 (V, p.~228) and 6 (V, p.~230) do not have converses: the existence of a comparison relation \(\mathbf{f} \preccurlyeq g\) , \(\mathbf{f} \ll g\) , \(\mathbf{f} \sim \mathbf{c}g\) between two functions that are differentiable on a neighbourhood of \(+\infty\) does not imply the same comparison relation between their derivatives, even for real monotone functions \(f\) and \(g\) .

For example, the function \(x^{2} + x\sin x + \cos x\) is monotone and equivalent to \(x^{2}\) , but its derivative \(x(2 + \cos x)\) is not equivalent to \(2x\) .

On the other hand, one can derive comparison relations when one assumes a priori that the derivatives of the functions considered are comparable (V, p.~217). In general, we shall say that two real functions f and g defined on an interval \([a, +\infty]\) are comparable to order k on a neighbourhood of \(+\infty\) if, on a neighbourhood of \(+\infty\) , they each have a regulated \(k^{th}\) derivative except at a countable set of points, and if, on this neighbourhood, \(f^{(k)}\) and \(g^{(k)}\) have constant sign (on the set on which they are defined) and are comparable.

We agree to say that two comparable real functions (V, p.~217) are comparable of order 0.

PROPOSITION 7. If two real functions \(f, g\) , are comparable of order 1, then they are comparable; further, the relation \(f \ll g\) (resp. \(f \sim cg\) , \(c\) constant) implies \(f' \ll g'\) (resp. \(f' \sim cg'\) ) unless \(g\) is equivalent to a nonzero constant.

Now, since \(f'\) and \(g'\) are of constant sign on an interval \([x_0, +\infty[\) , both \(f\) and \(g\) are monotone on this interval, so tend to a finite or infinite limit as \(x\) tends to \(+\infty\) . It is clear that \(f\) and \(g\) are comparable as \(x\) tends to \(+\infty\) if one of these limits is finite and \(\neq 0\) , or if one is zero and the other infinite. If both \(f\) and \(g\) tend to 0 one can write \(f(x) = -\int_x^{+\infty} f'(t) dt\) , \(g(x) = -\int_x^{+\infty} g'(t) dt\) ; since \(f'\) and \(g'\) are comparable the same is true for \(f\) and \(g\) and the comparison relation between \(f\) and \(g\) is the same as that between \(f'\) and \(g'\) , by prop. 6 (V, p.~20). Similarly, if \(f\) and \(g\) both have an infinite limit one has \(f(x) = f(x_0) + \int_{x_0}^x f'(t) dt\) , \(g(x) = g(x_0) + \int_{x_0}^x g'(t) dt\) ; again prop. 6 (V, p.~230) shows that \(f\) and \(g\) are comparable and that the comparison relation between \(f\) and \(g\) is the same as that between \(f'\) and \(g'\) . To complete the proof it remains to consider the case where \(g\) tends to \(\pm \infty\) and \(f\) to a constant; one cannot then have \(f' \succcurlyeq g'\) , since then one could deduce from prop. 1 (V, p.~218) that the integral \(\int_{x_0}^{+\infty} g'(t) dt\) was convergent; since \(f'\) and \(g'\) have been assumed comparable, one must have \(f' \ll g'\) .

COROLLARY. If two real functions \(f, g\) , are comparable of order \(k \geqslant 1\) , then they are comparable of order \(p\) for \(0 \leqslant p \leqslant k\) ; further, the relation \(f \ll g\) (resp. \(f \sim cg\) ) implies \(f^{(k)} \ll g^{(k)}\) (resp. \(f^{(k)} \sim cg^{(k)}\) ) unless one of the derivatives \(g^{(p)}\) ( \(0 \leqslant p \leqslant k - 1\) ) is equivalent to a constant \(\neq 0\) .

Indeed, since \(f^{(k)}\) and \(g^{(k)}\) have constant sign on an interval \([x_0, +\infty[\) , it follows that \(f^{(k-1)}\) and \(g^{(k-1)}\) are monotone on this interval, so have constant sign on a neighbourhood of \(+\infty\) ; further, prop. 7 of V, p.~232, shows that \(f^{(k-1)}\) and \(g^{(k-1)}\) are comparable, so the corollary follows from applying prop. 7 recursively.

Remarks. 1) The restriction on \(g\) in prop. 7 is essential. For example, one has \(\frac{1}{x} \ll 1 + \frac{1}{x}\) even though the derivatives of the two sides are equivalent; similarly \(1 + \frac{1}{x} \sim 1 + \frac{1}{x^2}\) , but \(1 / x^2 \gg 2 / x^3\) .

\begin{enumerate}
\def\labelenumi{\arabic{enumi})}
\setcounter{enumi}{1}
\item
  Though \(f\) and \(g\) are comparable of order \(k\) a function \(f_{1}\) equivalent to \(f\) need not be comparable of order \(k\) to \(g\) ; however it will be so if one assumes that \(f_{1}\) is comparable of order \(k\) to \(f\) and that none of the derivatives \(f^{(p)}\) ( \(0 \leqslant p \leqslant k - 1\) ) is equivalent to a nonzero constant.
\item
  If f and g are comparable of order k this need not be so for hf and hg even for a monotone function h as simple as \(h(x) = x\) (V, p.~260, exerc. 3); likewise, 1/f and 1/g are not necessarily comparable of order k (V, p.~259, exerc. 1).
\end{enumerate}

